% ── Section IV: Experimental Setup ───────────────────────────────────
% Conference budget: ~1 page.
% NOTE: ablation order fixed to match the locked scope + configs/:
%   Exp1 = +C2 (CTRV), Exp2 = C1+C2, Exp3 = C1+C2+C3.
%   (An earlier skeleton had Exp1 = +C1 — wrong.)
% ─────────────────────────────────────────────────────────────────────

\section{Experimental Setup}
\label{sec:experiments}

\subsection{Datasets}
\label{subsec:datasets}

\paragraph{VN-MixedTraffic (primary, self-collected)}
Fixed surveillance-style cameras over Vietnamese urban intersections
and arterials with motorcycle-dominant mixed traffic.
Annotated classes: \texttt{car}, \texttt{motorcycle}, \texttt{bus},
\texttt{truck}, \texttt{accident}. Any future accident-event evaluation requires
positive/negative clip labels and event onset/end annotations.
\todo{Fill after dataset freeze: number of clips, total duration,
resolution/fps, class-distribution table, accident/non-accident clip counts,
train/test protocol.}
The dataset is not publicly released for intellectual-property
reasons; evaluation protocol and pseudocode are fully specified.

\paragraph{Tracking ground truth and its provenance}
Tracking metrics are computed on a $600$-frame segment
($800\times450$, $\sim$47 tracks/frame), chosen as the densest window of
the sequence and lying in a held-out region.
It contains $27\,513$ boxes over $335$ identities, annotated in CVAT
under a fixed convention: a motorcycle box covers rider and passengers;
objects below $\sim$8\,px are not labeled; and a fully occluded track is
allowed to break but must be re-joined to its original identity on
reappearance, so that re-identification remains measurable.

Annotators corrected \emph{pre-labels} rather than drawing from scratch,
and those pre-labels came from the proposed tracker.
We state this because it biases the proposed method's own row: an error
the annotator does not catch survives as agreement.
The correction was substantial ($873$ boxes added, $1388$ deleted,
$1019$ adjusted, $15$ merges, $9$ splits --- $\sim$8\% of boxes
touched), and we bound the residual bias directly in
Sec.~\ref{subsec:ablation_results}.
We verify the annotation pipeline with three checks: the ground truth
scored against itself gives MOTA${=}$IDF1${=}1.000$; frame alignment
peaks unambiguously at zero offset (mean best IoU $0.960$, against
$0.705$ and $0.715$ at $\pm1$); and the added/deleted box counts match
the seed model's miss/false-positive counts exactly, closing the
export--annotate--import loop with no free parameters.

\paragraph{Public benchmarks}
CADP~\cite{Shah2018_cadp} (CCTV accident clips) and
D2-City~\cite{Che2019_d2city} (large-scale dashcam footage) are used
for comparison with prior accident-detection results only --- no
training is performed on them.

\subsection{Metrics}
\label{subsec:metrics}

\emph{Tracking:} MOTA~\cite{Bernardin2008_mota}, IDF1, identity
switches (IDSW), and track fragmentation rate.
Our metric implementation is written from scratch and cross-validated
against \texttt{py-motmetrics} on shared sequences.
\emph{Detection:} precision, recall, F1, and mAP@0.5.
\emph{Accident branch (planned, untested):} binary event-level alert
precision/recall/F1, false alerts per hour of video and detection delay.
No accident-alert performance is established by the current tracking evaluation.
\emph{System:} per-frame latency (ms) and end-to-end FPS.

\subsection{Implementation Details}
\label{subsec:impl}

The detector is YOLO11~\cite{Khanam2024_yolo11}
(\todo{variant, training epochs, image size, batch size}) fine-tuned
on VN-MixedTraffic. A binary event-alert implementation and evaluation remain
future work; no downstream classifier training is claimed here.
The tracker is ByteTrack~\cite{Zhang2022_bytetrack} with the motion
model made class-conditional (Sec.~\ref{subsec:c2}); association logic
and both association stages are unchanged.
CTRV uses $\Delta t = 1$ frame with process noise
$\mathbf{Q} = \mathrm{diag}(0.5, 0.5, 0.1, 10^{-3}, 10^{-4})$ and
measurement noise $\mathbf{R} = \mathrm{diag}(10, 10)$ (pixels), tuned
on synthetic trajectories and left unchanged for the real-footage
evaluation; NSA scales $\mathbf{R}$ per detection at run time
(Eq.~\ref{eq:nsa}).
C1 uses $\alpha = 0.1$ and hysteresis $N_h = 15$ frames; its
\textsc{dense-peak} row is calibrated by grid search against the
annotated segment (Table~\ref{tab:param_table}), and the remaining rows
keep their initial values because no annotated footage exercises them.
C3 uses the weights and thresholds of Sec.~\ref{subsec:c3}.
Software: Python 3.10, PyTorch~2.x.
Latency is measured on CPU (Intel Core i5-12500H, 16\,GB RAM) with the
detector excluded, so the reported figures isolate tracker-side cost and
are not an end-to-end throughput claim.

\subsection{Ablation Design}
\label{subsec:ablation_design}

Table~\ref{tab:ablation_design} lists the four configurations; each is
evaluated on the same VN-MixedTraffic test split with identical
detections wherever the configuration permits, isolating each
component's contribution.

\begin{table}[h]
  \centering
  \caption{Ablation configurations. \checkmark~=~component active.}
  \label{tab:ablation_design}
  \begin{tabular}{lcccl}
    \toprule
    Configuration & C1 & C2 & C3 & Description \\
    \midrule
    Exp~0 (baseline) & --         & --         & --         & YOLO11 + ByteTrack \\
    Exp~1            & --         & \checkmark & --         & + CTRV motion model \\
    Exp~2            & \checkmark & \checkmark & --         & + context adaptation \\
    Exp~3 (proposed) & \checkmark & \checkmark & \checkmark & + reliability score \\
    \bottomrule
  \end{tabular}
\end{table}

\subsection{Baselines}
\label{subsec:baselines}

(i)~\emph{Vanilla ByteTrack} --- the direct baseline, differing only in
the proposed components.
(ii)~\emph{YOLO-only alerting} (no tracker) --- quantifies how much
temporal confirmation alone reduces false alerts.
(iii)~\emph{Hikvision AID}~\cite{Hikvision_AID} --- commercial
reference, cited as prior art; since its internals and operating points
are unpublished, the comparison is qualitative (capability matrix), not
numeric.
